\section{Marchenko--Pastur 律给出谱的精确形状}
\label{sec:11-Mathematical-Statistics/09-Random-Matrix-Theory/02}

你在特征平台上做维度诊断：\(p=400\) 个行为特征，\(n=2000\) 条样本，标准化之后算相关矩阵，特征值从 8.4 一路降到 0.02。碎石图上前面十几个明显高出一截，然后是一段说不清是不是肘部的缓坡。领导问保留几个主成分，你只能说``看图大概十来个''。问题在于这句话没有对照组：如果这 400 个特征之间根本没有任何相关性，特征值本来该有多高？没有这个基准，肘部是拍出来的。

上一节说样本谱会被推开一片。这一节把那片区域的边界和形状写成公式，于是``纯噪声该长什么样''变成一个能算的数。

\subsection{纯噪声的谱有一条闭式密度曲线}

设 \(X_1,\ldots,X_n\in\R^p\) 独立同分布、均值为零、协方差为 \(I\)、四阶矩有限，样本协方差

\[
S_n=\frac1n\sum_{i=1}^n X_iX_i^{\trans}.
\]

当 \(p,n\to\infty\) 且 \(p/n\to c\in(0,1]\) 时，\(S_n\) 的经验谱分布——把 \(p\) 个特征值当成 \(p\) 个点画出的直方图——几乎必然收敛到 Marchenko--Pastur 分布，它在区间 \([a,b]\) 上有密度

\[
f_c(x)=\frac{\sqrt{(b-x)(x-a)}}{2\pi cx},
\qquad a=(1-\sqrt c)^2,\quad b=(1+\sqrt c)^2 .
\]

\(c>1\) 时秩不足，谱在 0 处带一块质量 \(1-1/c\)，其余质量 \(1/c\) 仍按上式的形状铺在 \([a,b]\) 上。

三个读数值得记住。支撑区间的宽度是 \(b-a=4\sqrt c\)，随 \(c\) 以平方根速率张开；上界 \(b=(1+\sqrt c)^2\) 是最大特征值的落点；\(c\to0\) 时 \([a,b]\to\{1\}\)，低维直觉作为极限情形被包含在内。回到上一节的 \(c=1\)，公式给出 \([0,4]\)，与模拟里那条从 0 铺到 4 的直方图对上了。

\begin{center}
\begin{tikzpicture}[>=stealth,x=0.82cm,y=1cm]
\draw[->] (-0.2,0) -- (7.6,0) node[right,font=\small] {\(\lambda\)};
\draw[->] (0,-0.15) -- (0,3.1) node[above,font=\small] {密度};
\foreach \x/\l in {0/0,1/1,2/2,3/3,4/4,5/5,6/6,7/7}
  \draw (\x,0) -- (\x,-0.12) node[below,font=\scriptsize] {\l};
\fill[black!12] (0.26,0)
  -- plot[domain=0.26:2.24,samples=150,smooth] (\x,{2.6*sqrt((2.25-\x)*(\x-0.25))/(1.5708*\x)})
  -- (2.24,0) -- cycle;
\draw[thick,domain=0.26:2.24,samples=150,smooth]
  plot (\x,{2.6*sqrt((2.25-\x)*(\x-0.25))/(1.5708*\x)});
\draw[dashed] (2.25,0) -- (2.25,2.75);
\node[above,font=\scriptsize,anchor=south west] at (2.3,2.55) {上界 \(b=(1+\sqrt c)^2\)};
\node[font=\scriptsize,anchor=north] at (1.15,-0.45) {噪声区：区分不开};
\draw[very thick] (3.3,0) -- (3.3,1.25);
\draw[very thick] (4.6,0) -- (4.6,1.05);
\draw[very thick] (6.4,0) -- (6.4,0.85);
\fill (3.3,1.25) circle (2pt);
\fill (4.6,1.05) circle (2pt);
\fill (6.4,0.85) circle (2pt);
\node[font=\scriptsize,anchor=west] at (3.5,1.75) {越界的特征值：候选信号};
\draw[->] (4.3,1.62) -- (3.4,1.35);
\end{tikzpicture}
\end{center}

\emph{\(c=0.25\) 的 MP 密度铺在 \([0.25,2.25]\) 上；越过虚线的那几根竖线才有资格被当作信号。}

图里那条虚线是本节的产出。灰色区域里的特征值，无论看起来多高多低，都在纯噪声的正常范围内；只有右侧那几根穿过虚线的，才值得进一步追问。

\subsection{证明骨架与它依赖的条件}

完整证明不在这里展开，但骨架值得知道，因为条件都挂在骨架上。

工具是 Stieltjes 变换 \(s(z)=\int(x-z)^{-1}\,\mathrm d\mu(x)\)，它把测度编码成上半平面上的解析函数，而 \(\frac1\pi\operatorname{Im}s(x+\mathrm i0^+)\) 把密度还原回来。第一步是截断与重新中心化，把矩条件转化为对元素大小的一致控制。第二步是对 \(s_n(z)=\frac1p\trace(S_n-zI)^{-1}\) 用降秩恒等式：抽掉一个样本，用 Schur 补把 \(p\times p\) 的逆的对角元写成一个二次型，这个二次型在独立性下集中到它的期望附近。第三步取极限，集中后的递推退化成一个二次方程

\[
czs^2-(1-c-z)s+1=0 ,
\]

它的判别式恰好是 \((z-a)(z-b)\)——支撑区间的两个端点就是这样冒出来的，不是凑出来的。在上半平面解析且满足 \(s(z)\sim-1/z\) 的那一支是唯一的，取虚部即得 \(f_c\)。

条件对账：四阶矩有限用在第一步的截断上，重尾数据这一步就过不去；独立同分布用在第二步的集中上，样本之间有相关（时间序列的自相关是最常见的一种）会让二次型不再集中；\(p/n\to c\) 用在第三步的取极限上，它不要求 \(n\) 多大，只要求比值稳定。至于数据的具体分布，骨架里从头到尾没用上——Gauss、Bernoulli、均匀分布给出同一条曲线。这条普适性让 MP 律可以从无线通信的天线信号直接搬到金融的收益率矩阵，也是它能当通用基准的原因。

\subsection{把上界当成信号阈值，是这条定律最直接的用法}

现在回到开头那 400 个特征。\(c=400/2000=0.2\)，\(\sqrt c=0.447\)，上界 \(b=(1+0.447)^2=2.09\)。于是有了一条可执行的规则：数一数有几个特征值超过 2.09。假设是 6 个，那么剩下 394 个方向上的``结构''在统计上与纯噪声无法区分，碎石图上那十来个高点里有一多半是噪声的正常起伏。保留 6 个主成分有据可依，保留 12 个是在给噪声建模。

这条规则背后是一个假设检验的形状。零假设是``特征之间无相关''，检验统计量是最大特征值，而 MP 律给出了它在零假设下的落点。落点还很锐利：\(\lambda_{\max}\) 在 \(b\) 附近的涨落是 \(n^{-2/3}\) 量级，服从 Tracy--Widom 分布，比通常的 \(n^{-1/2}\) 还小一个数量级。这意味着阈值不是一条模糊的带，而是一条相当细的线；真要控制第一类错误，可以按 Tracy--Widom 的分位数在 \(b\) 上加一个 \(n^{-2/3}\) 量级的余量，但对 \(n\) 上千的数据，这个余量通常小到不影响计数。

用起来还有两处细节。其一，MP 律是按``真协方差为 \(I\)''陈述的，所以要先把每个特征标准化到方差 1，实际上是在对相关矩阵而不是原始协方差矩阵用阈值；量纲不同的特征不做标准化，最大特征值只是在报告哪个特征的单位最大。其二，若各特征方差不等但仍无相关，可以先用 \(\trace S_n/p\) 估出噪声尺度 \(\sigma^2\)，把阈值放到 \(\sigma^2(1+\sqrt c)^2\) 上。金融领域早就在这么做：把数百只股票的日收益率相关矩阵拿来数越界特征值，越界的往往只有百分之几，对应市场因子和几个行业因子，其余绝大多数特征值老老实实落在 MP 区间里——那些方向上看似存在的``板块结构''，与随机数据给出的结构没有区别。

\subsection{阈值失效的几种方式都指向同一句话：先检查零假设}

MP 律作为基准的前提，是那些条件确实成立。它们不成立时，越界的特征值可能与信号无关。

样本之间相关是最常见的破坏方式。日频收益率、传感器时序、同一用户的多条行为记录，行与行之间不独立，有效样本量远小于名义的 \(n\)。此时真实的 \(c\) 比 \(p/n\) 大，真实上界比 \(2.09\) 高，用 \(2.09\) 去数就会把噪声数成信号。补救是按块或按时间窗估计有效样本量，或者干脆用置换检验：打乱每一列的行顺序、破坏跨特征相关但保留每列的边缘分布，重算最大特征值，重复几百次得到经验零分布。这个办法不依赖矩条件，代价是算力。

特征之间有确定性依赖是第二种。同一个类别变量展开的独热编码之和恒为 1，重复采集的特征几乎共线，这些结构会造出与样本量无关的大特征值。它们越界，但越的不是信号的界，而是你在特征工程里自己埋进去的界。

重尾是第三种。四阶矩发散时最大特征值不再收敛到 \(b\)，而是由个别极端观测主导，越界几乎必然发生。先做稳健标准化或秩变换，比事后解释那个大特征值划算。

\(\Sigma\ne I\) 则不算失效，只算换了一条曲线。此时极限谱由 \(\Sigma\) 的谱与 MP 律的自由卷积给出，形状更复杂但仍然确定；这也正是特征值收缩类方法的出发点——既然知道从真谱到样本谱的``涂抹''规则，就可以反过来把样本特征值往回推。Ledoit--Wolf 收缩把 \(\widehat\Sigma\) 向 \(\bar\lambda I\) 拉，收缩强度的最优取值随 \(c\) 增大而增大，这与 \(4\sqrt c\) 的弥散宽度是同一个量在两处的表现；\hyperref[sec:13-Machine-Learning/06-Probabilistic-Models/02]{``LDA 与收缩：判别方向的几何与高维失效''}和\hyperref[sec:09-Convex-Optimization/07-Applications/04]{``风险、收益与投资组合''}分别在分类和组合优化里用了这件事。

到这里，噪声的形状已经有了公式，信号的门槛也有了数值。还剩一个问题没答：一个特征值刚刚越过 \(b\)，它对应的特征向量能信吗？下一节说明，这两件事的答案并不同步。
